\textbf{1.} Vrai, et même sans seuil : tout entier $n \geqslant 3$ s'écrit
$n = 3 + (n - 3)$, et tout entier $n$ s'écrit $n = 0 + n$, où $0$ est un
multiple de $3$. Un énoncé vrai se \emph{démontre} : ici, une écriture
explicite valable pour tout $n$ suffit. (Ne pas confondre avec la conjecture
de Goldbach, \og tout entier pair $n \geqslant 4$ est somme de deux nombres
premiers \fg, qui n'est toujours pas démontrée.)

\textbf{2.} Faux : pour $n = 2$, $3^4 + 1 = 82$ n'est pas divisible par $10$.
(L'énoncé est vrai pour $n = 1$ : $3^2 + 1 = 10$. On peut montrer qu'il est
vrai exactement pour les $n$ impairs, car $3^{2n} = 9^n$ se termine par $9$
si $n$ est impair et par $1$ si $n$ est pair.)

\textbf{3.} Vrai, pour tout polygone à $n \geqslant 3$ côtés, convexe ou non.
Une diagonale est un segment joignant deux sommets non consécutifs. De
chaque sommet partent $n - 3$ diagonales (vers tous les sommets sauf
lui-même et ses deux voisins). Cela fait $n(n-3)$ extrémités de diagonales,
et chaque diagonale est comptée deux fois (une fois par extrémité) : il y a
$\frac{n(n-3)}{2}$ diagonales. Le raisonnement n'utilise ni la régularité ni
la convexité : la forme du polygone n'intervient pas.

\textbf{4.} Faux : $4 + 8 = 12$ n'est pas un multiple de $8$.

\textbf{5.} Faux : un losange dont les angles ne sont pas droits (par exemple
de sommets $(0\,;0)$, $(2\,;1)$, $(4\,;0)$, $(2\,;-1)$, de côtés tous égaux à
$\sqrt5$) n'est pas un carré. \hfill $\blacksquare$
